\PuzzleChapterIntro{The Twin-Wick Corridor}{Lore Logic \textbullet\ Counting by Parity \textbullet\ Unique Integer}{This problem asks for a ratio, which suggests you should not compute the massive numerator and denominator separately. Instead, look for a structural relationship: given a specific sequence of stations, exactly how many ways can you assign the wicks?}
\Lore{
The Corridor of Wicks is lit in pairs: each station a matched cradle of brass with a single number stamped deep enough to catch soot.

At every station a Sun wick and a Shade wick share that number; a pull flips the state with a soft click, lit to dark and dark to lit.

The steward keeps the corridor quiet, because the sound of toggling carries.

An old line is etched into the iron at eye height:
\begin{quote}
\textit{``The Shade must never shine---unless you walk the reckless road.''}
\end{quote}

Two rites are copied in the ledger. The Vault accepts only their ratio.
}

\Problem

There are \textbf{9 numbered stations}, labeled \(1,2,\dots,9\).
At each station \(i\) there are exactly two wicks: a \textbf{Sun} wick and a \textbf{Shade} wick.
All wicks start \textbf{dark}.

A \textbf{pull} chooses a station \(i\) and then chooses \emph{exactly one} of its two wicks; that wick flips state (dark \(\leftrightarrow\) lit).

A \textbf{rite} is a sequence of exactly \textbf{21 pulls}.
Two rites are distinct iff they differ at some pull index.

At the end of the 21 pulls, the door demands:
\begin{itemize}
  \item all 9 Sun wicks are \textbf{lit}, and
  \item all 9 Shade wicks are \textbf{dark}.
\end{itemize}

Define:
\begin{itemize}
  \item \(N\): the number of \textbf{Reckless Rites} that satisfy the final demand (Shade wicks may be lit during the rite).
  \item \(M\): the number of \textbf{Oathbound Rites} that satisfy the final demand and in which
  \textbf{after every prefix of the 21 pulls, all 9 Shade wicks are dark}.
\end{itemize}

\VaultDirective{Compute the integer factor $\dfrac{N}{M}$.}

\SolutionPage

A wick ends \textbf{lit} exactly when it is pulled an \emph{odd} number of times, and ends \textbf{dark} exactly when it is pulled an \emph{even} number of times.

\medskip
\textbf{ Oathbound rites never pull Shade.}

In an Oathbound rite, each Shade wick is dark after every prefix.
Since every Shade wick starts dark, pulling a Shade wick even once would make it lit immediately after that pull, violating the oath.
So in an Oathbound rite, \emph{no Shade wick is ever pulled}.

Therefore an Oathbound rite consists of 21 pulls on Sun wicks only.
To satisfy the final demand, each Sun wick must be pulled an odd number of times.
So \(M\) counts exactly the length-21 sequences of Sun-pulls in which each of the 9 Sun wicks is used an odd number of times.

\medskip
\textbf{The station-number skeleton.}

Given any rite, forget whether each pull was Sun or Shade and record only its station number.
This produces a length-21 \emph{skeleton} word in \(\{1,2,\dots,9\}\).

For a successful Reckless rite, fix a station \(i\) and write
\[
s_i=\#(\text{Sun pulls at station }i),\qquad h_i=\#(\text{Shade pulls at station }i),\qquad c_i=s_i+h_i.
\]
The final demand requires \(s_i\) to be odd (Sun ends lit) and \(h_i\) to be even (Shade ends dark), hence \(c_i\) is odd.
So in the skeleton, each symbol \(i\) appears an odd number of times.

Conversely, any skeleton in which each \(i\in\{1,\dots,9\}\) appears an odd number of times determines exactly one Oathbound rite:
at each occurrence of \(i\), pull the \emph{Sun} wick at station \(i\).
This is valid and satisfies the final demand.

Hence the set of possible skeletons is in bijection with the set of Oathbound rites, so the number of skeletons is exactly \(M\).

\medskip
\textbf{Expanding one skeleton into Reckless rites.}

Fix one skeleton, and let \(c_i\) be the number of occurrences of station \(i\) in it.
Then each \(c_i\) is odd (hence \(c_i\ge 1\)) and \(\sum_{i=1}^9 c_i=21\).

To obtain a successful Reckless rite from this skeleton, for each station \(i\) we must decide, among the \(c_i\) occurrences of \(i\),
which are \emph{Shade} pulls (the rest are \emph{Sun} pulls).
The final demand requires the Shade wick to end dark, so the number of Shade pulls at station \(i\) must be even.
Thus we must choose an even-sized subset of the \(c_i\) occurrences.

Among the \(2^{c_i}\) subsets, exactly half have even size (toggle a distinguished occurrence to pair even and odd subsets),
so there are \(2^{c_i-1}\) valid choices for station \(i\).

These choices are independent across stations, so the number of successful Reckless rites with this fixed skeleton is
\[
\prod_{i=1}^9 2^{c_i-1}
=2^{(c_1+\cdots+c_9)-9}
=2^{21-9}
=2^{12}
=4096.
\]

Thus each skeleton contributes exactly \(4096\) Reckless rites. Since the number of skeletons is \(M\), we have \(N=4096\,M\), and therefore
\[
\frac{N}{M}=4096.
\]

\FinalAnswer{4096}

\NextPage
